\documentclass[12pt,a4paper]{article}
\usepackage[spanish]{babel}
\usepackage[utf8]{inputenc}
\usepackage[T1]{fontenc}
\usepackage{geometry}
\usepackage{graphicx}
\usepackage{xcolor}
\usepackage{fancyvrb}
\usepackage{enumitem}
\geometry{margin=2.2cm}

\title{Reporte: 5 Operaciones con Git}
\author{Repositorio: formularioPaises}
\date{08 de abril de 2026}

\begin{document}
\maketitle

\section*{Objetivo}
Documentar 5 operaciones de Git con palabras simples y claras, incluyendo:
\begin{itemize}[leftmargin=*]
\item Comando o programa.
\item Explicacion de parametros.
\item Problema o situacion donde se aplico.
\item Captura de pantalla.
\item Efecto obtenido.
\end{itemize}

\section*{Nota importante sobre capturas}
En este documento incluyo dos tipos de captura:
\begin{itemize}[leftmargin=*]
\item \textbf{Captura textual}: salida real en bloque de terminal o vista de cambios.
\item \textbf{Espacio para imagen}: para que pegues tu screenshot si el profesor pide imagen visual.
\end{itemize}

\hrule
\vspace{0.4cm}

\section*{Operacion 1: Ver estado general}
\textbf{Comando o programa}
\begin{Verbatim}[fontsize=\small,frame=single]
git status --short
\end{Verbatim}

\textbf{Explicacion de los parametros}
\begin{itemize}[leftmargin=*]
\item \texttt{--short}: muestra el estado en formato corto (mas facil de leer).
\end{itemize}

\textbf{Problema o situacion donde se aplico}
\begin{itemize}[leftmargin=*]
\item Queria revisar rapido que archivos estaban cambiados despues de modificar la API.
\end{itemize}

\textbf{Captura de pantalla (texto)}
\begin{Verbatim}[fontsize=\small,frame=single]
Resumen observado en el workspace:
- Hay cambios sin stage en varios archivos.
- Ejemplos: python-api/app.py, python-api/README.md,
  y varios archivos dentro de formulario-paises-app/.
\end{Verbatim}

\textbf{Captura de pantalla (imagen, opcional)}
\begin{itemize}[leftmargin=*]
\item Guarda una imagen como \texttt{capturas/op1-status.png} y descomenta esta linea:
\end{itemize}
% \includegraphics[width=0.95\textwidth]{capturas/op1-status.png}

\textbf{Efecto obtenido}
\begin{itemize}[leftmargin=*]
\item Se confirmo que el repositorio tiene cambios pendientes por organizar.
\end{itemize}

\hrule
\vspace{0.4cm}

\section*{Operacion 2: Revisar diferencia del backend}
\textbf{Comando o programa}
\begin{Verbatim}[fontsize=\small,frame=single]
git diff -- python-api/app.py
\end{Verbatim}

\textbf{Explicacion de los parametros}
\begin{itemize}[leftmargin=*]
\item \texttt{diff}: compara cambios no confirmados.
\item \texttt{-- python-api/app.py}: limita la comparacion a un archivo.
\end{itemize}

\textbf{Problema o situacion donde se aplico}
\begin{itemize}[leftmargin=*]
\item Necesitaba validar que la API paso de archivo de texto a SQLite.
\end{itemize}

\textbf{Captura de pantalla (texto)}
\begin{Verbatim}[fontsize=\small,frame=single]
Fragmentos del cambio observado:
+ import sqlite3
+ ARCHIVO_BASE_DATOS = Path(__file__).with_name('organizaciones.db')
+ def inicializar_base_datos():
+   CREATE TABLE IF NOT EXISTS organizaciones (...)
+ INSERT INTO organizaciones (...)
\end{Verbatim}

\textbf{Captura de pantalla (imagen, opcional)}
% \includegraphics[width=0.95\textwidth]{capturas/op2-diff-app.png}

\textbf{Efecto obtenido}
\begin{itemize}[leftmargin=*]
\item Se comprobo que el backend ahora guarda en base de datos SQLite.
\end{itemize}

\hrule
\vspace{0.4cm}

\section*{Operacion 3: Revisar diferencia de documentacion}
\textbf{Comando o programa}
\begin{Verbatim}[fontsize=\small,frame=single]
git diff -- python-api/README.md
\end{Verbatim}

\textbf{Explicacion de los parametros}
\begin{itemize}[leftmargin=*]
\item \texttt{diff}: muestra diferencia de contenido.
\item \texttt{-- python-api/README.md}: revisa solo el README de la API.
\end{itemize}

\textbf{Problema o situacion donde se aplico}
\begin{itemize}[leftmargin=*]
\item Queria asegurar que la documentacion coincida con el nuevo guardado en SQLite.
\end{itemize}

\textbf{Captura de pantalla (texto)}
\begin{Verbatim}[fontsize=\small,frame=single]
Cambio principal observado:
- guardar organizaciones en archivo de texto
+ guardar organizaciones en base de datos SQLite
- organizaciones.txt
+ organizaciones.db
\end{Verbatim}

\textbf{Captura de pantalla (imagen, opcional)}
% \includegraphics[width=0.95\textwidth]{capturas/op3-diff-readme.png}

\textbf{Efecto obtenido}
\begin{itemize}[leftmargin=*]
\item La documentacion quedo coherente con el codigo real.
\end{itemize}

\hrule
\vspace{0.4cm}

\section*{Operacion 4: Preparar archivos para commit}
\textbf{Comando o programa}
\begin{Verbatim}[fontsize=\small,frame=single]
git add python-api/app.py python-api/README.md
\end{Verbatim}

\textbf{Explicacion de los parametros}
\begin{itemize}[leftmargin=*]
\item \texttt{add}: pasa cambios al area de stage.
\item \texttt{python-api/app.py python-api/README.md}: solo agrega esos 2 archivos.
\end{itemize}

\textbf{Problema o situacion donde se aplico}
\begin{itemize}[leftmargin=*]
\item Habia muchos cambios en el repo y se queria confirmar solo lo del backend.
\end{itemize}

\textbf{Captura de pantalla (texto sugerido para tu terminal)}
\begin{Verbatim}[fontsize=\small,frame=single]
Despues de ejecutar:
git status --short

Esperado:
A  python-api/app.py
A  python-api/README.md
(El resto sigue como cambios sin stage)
\end{Verbatim}

\textbf{Captura de pantalla (imagen, opcional)}
% \includegraphics[width=0.95\textwidth]{capturas/op4-add.png}

\textbf{Efecto obtenido}
\begin{itemize}[leftmargin=*]
\item Se separan los cambios importantes del backend del resto de archivos.
\end{itemize}

\hrule
\vspace{0.4cm}

\section*{Operacion 5: Verificar stage antes de commit}
\textbf{Comando o programa}
\begin{Verbatim}[fontsize=\small,frame=single]
git diff --cached --name-only
\end{Verbatim}

\textbf{Explicacion de los parametros}
\begin{itemize}[leftmargin=*]
\item \texttt{diff --cached}: revisa lo que esta en stage.
\item \texttt{--name-only}: muestra solo nombres de archivo.
\end{itemize}

\textbf{Problema o situacion donde se aplico}
\begin{itemize}[leftmargin=*]
\item Antes de commit, queria verificar exactamente que se iba a confirmar.
\end{itemize}

\textbf{Captura de pantalla (texto)}
\begin{Verbatim}[fontsize=\small,frame=single]
Si ya hiciste git add de la Operacion 4, debe salir:
python-api/app.py
python-api/README.md
\end{Verbatim}

\textbf{Captura de pantalla (imagen, opcional)}
% \includegraphics[width=0.95\textwidth]{capturas/op5-cached.png}

\textbf{Efecto obtenido}
\begin{itemize}[leftmargin=*]
\item Se evita hacer un commit con archivos que no querias incluir.
\end{itemize}

\hrule
\vspace{0.4cm}

\section*{Cierre}
Con estas 5 operaciones se logro:
\begin{itemize}[leftmargin=*]
\item Revisar el estado general del repositorio.
\item Validar cambios clave en backend y README.
\item Preparar y verificar archivos para un commit limpio.
\end{itemize}

\section*{Comando final recomendado}
\begin{Verbatim}[fontsize=\small,frame=single]
git commit -m "feat(api): guardar organizaciones en SQLite"
\end{Verbatim}

\end{document}
